\documentclass[12pt]{article}
\usepackage[utf8]{inputenc}
\usepackage[T1]{fontenc}
\usepackage{hyperref}
\title{Kamang Language Resources for K-12 Curriculum Development}
\author{Compiled by Claude Sonnet 5.5 (claude-sonnet-5-5) for Andr\'as Kornai}
\date{}
\begin{document}
\maketitle

\section{Introduction}
Kamang (also called Woisika; ISO 639-3: woi; Glottocode: kama1365) is a Timor-Alor-Pantar language of Alor Island, Indonesia, with about 6,000 speakers. It is endangered, because children often have only passive competence and are shifting to Malay. The main dialects are Western, Lowland and Upland Kamang. It is not known to be taught in schools.

\section{Primary Educational Resources}
No primers, textbooks or school programmes were found.

\section{Linguistic and Reference Materials}
\begin{itemize}
\item Antoinette Schapper, ``Kamang'', in \emph{The Papuan Languages of Timor, Alor and Pantar, Volume 1: Sketch Grammars} (Mouton de Gruyter, 2014). It is the main grammar, and Kamang has more morphology than any other language in the family. Catalogue record: \url{https://oeaw.ac.at/resources/Record/9781614515241}
\item New Guinea World entry on Kamang. It includes pronoun and phonology data and cites the Schapper and Manimau (2011) Kamang dictionary, which includes some Takailubui, Suboo and South West variants. I saw no publisher or access link for the dictionary. \url{https://newguineaworld.linguistik.uzh.ch/families/trans-new-guinea/berau-gulf/west-bomberai/timor-alor-pantar/alor-pantar/east-alor/central-east-alor/kamang}
\item Schapper, ``Temperature terms in Kamang and Abui, two Papuan languages of Alor'' (Benjamins): \url{https://www.benjamins.com/catalog/tsl.107.29sch}
\item Schapper's author page listing her books: \url{https://cnrs.academia.edu/AntoinetteSchapper}
\item Wikipedia article on the language (as Woisika): \url{https://en.wikipedia.org/wiki/Woisika_language}
\end{itemize}

\section{Religious and Community Resources}
None found.

\section{Digital and Audio Resources}
None found.

\section{Curriculum Development Context}
The language is endangered and the documentation is mostly academic: a sketch grammar and a dictionary. Nothing was found that is aimed at children or schools. A curriculum would need a community-agreed orthography and new primers. Because children are shifting to Malay, early-grade materials would matter most.

\end{document}